\documentclass[11pt,a4paper]{article}

% --- Packages ---
\usepackage[margin=.75in]{geometry}
\usepackage[utf8]{inputenc}
\usepackage[T1]{fontenc}
\usepackage{lmodern} % Fix for font not found error
\usepackage{titlesec}
\usepackage{enumitem}
\usepackage{hyperref}
\usepackage{xcolor}
\usepackage{fancyhdr}
\usepackage{comment}
\usepackage{cleveref}
\usepackage[bottom]{footmisc}
\usepackage{graphicx}
\usepackage{longtable}
\usepackage{booktabs}
\usepackage{array}
\usepackage{calc}
\usepackage{etoolbox}

% Standard Pandoc setup
\providecommand{\tightlist}{%
  \setlength{\itemsep}{0pt}\setlength{\parskip}{0pt}}

\crefformat{section}{\S#2#1#3}
\crefformat{subsection}{\S#2#1#3}
\crefformat{subsubsection}{\S#2#1#3}
\crefformat{figure}{#2Figure~#1#3}
\crefformat{footnote}{#2\footnotemark[#1]#3}

\linespread{0.9} % Riduce lo spazio tra le linee

% --- Colors ---
\definecolor{primary}{RGB}{0, 0, 0}
\definecolor{links}{HTML}{0000EE}
\definecolor{blue}{HTML}{26428B}

\hypersetup{
    colorlinks=true,
    linkcolor=links,
    urlcolor=links,
    citecolor=links,
}

% --- Section Formatting ---
\titleformat{\section}
  {\normalfont\Large\bfseries\color{blue}}{\thesection}{1em}{}
\titleformat{\subsection}
{\normalfont\large\bfseries\color{blue}}{\thesubsection}{1em}{}
\titleformat{\subsubsection}
{\normalfont\bfseries\color{blue}}{\thesubsubsection}{1em}{}

% --- Header and Footer ---
\pagestyle{fancy}
\fancyhead[L]{\small Giuseppe Capasso}
\fancyhead[R]{\small intro-to-linux}
\fancyfoot[C]{\thepage}
\renewcommand{\headrulewidth}{0.4pt}

% --- Custom Title Command ---
\newcommand{\proposaltitle}[1]{
    \begin{center}
        \vspace*{0.1cm}
        {\LARGE \bfseries \color{blue} #1} \\[1.5ex]
        {\large \textbf{Giuseppe Capasso}} \\[1ex]
                \small{
            \vspace{0.1cm}
            \textbf{Materia:} intro-to-linux\\
        }
            \end{center}
    \vspace{0.3cm}
    \hrule height 0.5pt
    \vspace{0.5cm}
}

\begin{document}

\proposaltitle{Introduzione all'Ambiente Linux}

Introduzione all'Ambiente Linux Modulo: Sistemi Operativi Linux --
Lezione 1 Obiettivo: Analisi dell'architettura Linux, configurazione
dell'ambiente di virtualizzazione e interazione con la Shell.

\begin{enumerate}
\def\labelenumi{\arabic{enumi}.}
\tightlist
\item
  Architettura del Sistema Quando si parla di ``Linux'' in ambito
  tecnico, è necessario fare una distinzione tra il Kernel e il sistema
  operativo completo.
\end{enumerate}

1.1 Il Kernel Linux Linux è un Kernel monolitico open source (GPLv2),
creato inizialmente da Linus Torvalds nel 1991. Il Kernel è il
componente centrale del sistema operativo che risiede in memoria ed è
responsabile dell'interfaccia tra l'hardware della macchina e il
software applicativo. Le sue funzioni primarie includono: • Process
Management: Scheduling dei processi e gestione del multitasking. •
Memory Management: Gestione della RAM, paginazione e memoria virtuale. •
Hardware Abstraction: Fornitura di driver per interagire con le
periferiche (I/O). • System Calls: Interfaccia che permette ai programmi
in User Space di richiedere servizi al Kernel.

1.2 GNU/Linux (User Space) Un Kernel da solo non è utilizzabile
dall'utente finale. Il sistema operativo completo è costituito
dall'unione del Kernel Linux con il software di sistema del progetto GNU
(GNU's Not Unix), che fornisce librerie, compilatori e la shell.

\begin{enumerate}
\def\labelenumi{\arabic{enumi}.}
\setcounter{enumi}{1}
\tightlist
\item
  Ecosistema delle Distribuzioni Una Distribuzione Linux (Distro) è un
  sistema operativo completo che integra il Kernel, il set di strumenti
  GNU, un sistema di gestione dei pacchetti (Package Manager) e, spesso,
  un Desktop Environment. Le distribuzioni si differenziano
  principalmente per:
\end{enumerate}

1

1. Gestione dei Pacchetti: Il formato dei file di installazione e il
software che risolve le dipendenze. 2. Ciclo di Rilascio: Rolling
release (aggiornamenti continui) vs Fixed release (versioni stabili
periodiche).

Famiglie Principali • Debian/Ubuntu: Utilizza il formato .deb e il
package manager apt. Nota per la stabilità e l'ampia adozione sia
desktop che server. È la distribuzione di riferimento per questo corso.
• RHEL (Red Hat Enterprise Linux) / Fedora: Utilizza .rpm e dnf/yum. È
lo standard de facto per le infrastrutture enterprise aziendali. • Arch
Linux: Adotta un approccio ``minimalista'' e rolling release. Richiede
una configurazione manuale approfondita.

\begin{enumerate}
\def\labelenumi{\arabic{enumi}.}
\setcounter{enumi}{2}
\tightlist
\item
  Virtualizzazione (Lab Setup) Per l'ambiente di laboratorio
  utilizzeremo una soluzione di Virtualizzazione di Tipo 2 (Hosted).
\end{enumerate}

3.1 Concetti Chiave • Host OS: Il sistema operativo fisico
(Windows/macOS) su cui è installato l'hardware. • Hypervisor: Il
software (es. VirtualBox, VMware Workstation) che astrae le risorse
hardware dell'Host per assegnarle alle macchine virtuali. • Guest OS: Il
sistema operativo virtualizzato (nel nostro caso, Ubuntu).

3.2 Vantaggi in ambiente di sviluppo L'uso di VM permette l'isolamento
dei processi (sandbox). Eventuali errori di configurazione critica o
compromissioni di sicurezza nel Guest non impattano l'Host. Inoltre, la
funzionalità di Snapshot permette di salvare lo stato della memoria e
del disco per ripristinare il sistema a una condizione precedente in
caso di errore.

\begin{enumerate}
\def\labelenumi{\arabic{enumi}.}
\setcounter{enumi}{3}
\tightlist
\item
  Interazione con la Shell (CLI) La CLI (Command Line Interface) è
  l'interfaccia primaria per l'amministrazione di sistema (SysAdmin). In
  ambiente Linux, la shell predefinita è solitamente Bash (Bourne Again
  SHell). La CLI offre vantaggi sostanziali rispetto alla GUI: •
  Granularità: Accesso diretto alle System Call e configurazioni di
  basso livello.
\end{enumerate}

• Automazione: Possibilità di scripting per task ripetitivi. • Remote
Management: Gestione di server remoti tramite protocollo SSH, essenziale
in architetture Cloud/Headless.

4.1 Navigazione del Filesystem Il filesystem Linux segue lo standard FHS
(Filesystem Hierarchy Standard) ed è strutturato ad albero rovesciato,
partendo dalla radice / (root). • Path Assoluto: Percorso completo dalla
root (es. /home/studente/progetti). • Path Relativo: Percorso riferito
alla directory corrente (es. ./progetti). Comandi di Navigazione Comando
Descrizione tecnica Sintassi comune Print Working Directory. Restituisce
il path assoluto pwd pwd della directory corrente. ls -la (mostra
dettagli e file ls List. Elenca il contenuto di una directory.
nascosti/dotfiles) cd /path/to/dir o cd .. Change Directory. Modifica la
working directory cd corrente. (parent directory)

4.2 Manipolazione File e Directory Comando Descrizione tecnica mkdir
Make Directory. Crea una nuova directory. touch Aggiorna il timestamp di
un file o ne crea uno vuoto se non esiste. Copy. Copia file o directory
(con flag -r per cp ricorsivo). Move. Sposta o rinomina file/directory
(inode mv operation). rm

Remove. Rimuove file (unlink). Non esiste cestino.

Sintassi comune mkdir nome\_dir touch file.txt cp source dest mv
old\_name new\_name rm file o rm -rf dir (force recursive)

\begin{enumerate}
\def\labelenumi{\arabic{enumi}.}
\setcounter{enumi}{4}
\tightlist
\item
  Risorse e Riferimenti Tecnici Per l'approfondimento e l'esercitazione
  pratica pre-lab: • Linux Journey (linuxjourney.com): Consultare la
  sezione Grasshopper per la sintassi dei comandi. • Bandit Wargame
  (overthewire.org): Esercizi di sicurezza e amministrazione via SSH. 3
\end{enumerate}

• ExplainShell (explainshell.com): Tool per l'analisi visiva degli
argomenti dei comandi (parsing delle man pages). • Man Pages: Il manuale
di sistema integrato. Utilizzare man {[}comando{]} (es. man ls) per la
documentazione ufficiale POSIX.



\end{document}
